% ====================================================================
% JMP_seminar_deck_r6.tex -- the 17 September seminar deck, R6 content.
%
% AUTHORITY
%   docs/normative/JMP_W1_fork_ruling_v1.md, Appendix A (Deputy R1-R6)
%   docs/normative/JMP_W1_BASELINE_F1_authorization_and_Cobs_ruling_v1.md, s5
%
% R2 retires every slide, figure, number and note built on W1-EA.  None of
% that material is carried here, repaired here, or restated here; the v4.1
% deck keeps it as a historical artifact and is not built for this seminar.
%
% Every numeral comes from % deck_numbers_r6.tex -- GENERATED by make_deck_numbers_r6.py.
% Do not edit by hand.  Every numeral on an R6 slide comes from here.
% Sources: S11 parameter tables (model of record);
%          positive_fit_diagnostics_v2 (G2-ADEQUATE statistics only);
%          welfare/baseline_f1_v1 (BASELINE-F-1, commit 6048c9f, verified b5550af).

% --- BASELINE-F-1 (E1): household EUR/month, unequivalised, dwt-weighted ---
\newcommand{\WFSingN}{1\,540}
\newcommand{\WFSingMean}{1\,435}
\newcommand{\WFSingMedian}{1\,320}
\newcommand{\WFSingGini}{0.238}
\newcommand{\WFSingWorkers}{1\,336}
\newcommand{\WFSingNonworkers}{204}
\newcommand{\WFSingCone}{0}
\newcommand{\WFSingCtwo}{0}
\newcommand{\WFSingCfive}{1.8e-15}
\newcommand{\WFCoupN}{2\,223}
\newcommand{\WFCoupMean}{2\,496}
\newcommand{\WFCoupMedian}{2\,356}
\newcommand{\WFCoupGini}{0.210}
\newcommand{\WFCoupWorkers}{2\,173}
\newcommand{\WFCoupNonworkers}{50}
\newcommand{\WFCoupCone}{0}
\newcommand{\WFCoupCtwo}{0}
\newcommand{\WFCoupCfive}{1.8e-15}

% --- positive fit: ONLY G2-ADEQUATE statistics (weighted, the binding rule) ---
\newcommand{\FitExtCF}{89.7}
\newcommand{\FitExtRatioCF}{0.17}
\newcommand{\FitExtSF}{86.8}
\newcommand{\FitExtRatioSF}{0.24}
\newcommand{\FitAdequateCount}{2}
\newcommand{\FitLimitedCount}{2}

% --- S11 objectives reproduced by the diagnostics evaluator ---
\newcommand{\ObjSingles}{6253.463}
\newcommand{\ObjCouples}{10283.034}
\newcommand{\FitNodes}{100}
\newcommand{\FitSims}{500}
\newcommand{\FitBootReps}{200}
\newcommand{\EvaluatorCommit}{55bb0d0}

% --- S11 model of record: estimated opportunity kernel ---
\newcommand{\SSingParams}{52}
\newcommand{\SSingActive}{41}
\newcommand{\SElevenAccessConst}{-3.174}
\newcommand{\SElevenAccessConstZ}{-7.8}
\newcommand{\SElevenAccessSur}{-1.442}
\newcommand{\SElevenAccessSurZ}{-6.1}
\newcommand{\SElevenHoursFThirtyFive}{2.066}
\newcommand{\SElevenHoursFThirtyFiveZ}{23.1}
\newcommand{\SElevenBetaC}{2.039}
\newcommand{\SElevenBetaCZ}{7.0}
\newcommand{\SElevenWageSigma}{0.382}
\newcommand{\SElevenWageSigmaZ}{28.7}
\newcommand{\SElevenWageEducH}{0.149}
\newcommand{\SElevenWageEducHZ}{4.8}
\newcommand{\SElevenWagePexp}{0.241}
\newcommand{\SElevenWagePexpZ}{2.8}
\newcommand{\SCoupParams}{47}
\newcommand{\SCoupActive}{47}

% --- estimation frames (S11 corrected criterion-A) ---
\newcommand{\FrameSingles}{1\,540}
\newcommand{\FrameCouples}{2\,223}

, generated by
% make_deck_numbers_r6.py from three authorized sources only.  No numeral is
% typed on a slide.
% ====================================================================
\documentclass[10pt,aspectratio=169]{beamer}
% ====================================================================
% jmp_beamer_preamble.tex --- shared preamble for the JMP seminar deck.
%
% REUSED VERBATIM from the companion theory talk
%   beamer/reference/Theory_talk/slides/Presentation.tex
%     "Jobs and Well-Being Measurement" (Haydar and Maniquet)
% so that the two decks read as one visual family.  The elements carried
% over unchanged are marked [THEORY-TALK] below; everything else is an
% addition this deck needs and the theory talk does not.
%
% Consistency is the point: same class options, same theme, same accent
% colour, same frame numbering, same appendix handling, same visible-link
% convention, same note-page template, same two-agent sketch palette.
% ====================================================================

% ---- [THEORY-TALK] class, theme and the deep-red accent -------------
\usetheme{metropolis}
\definecolor{deepred}{RGB}{150,25,35}
\setbeamercolor{frametitle}{bg=deepred}
\setbeamertemplate{navigation symbols}{}
% metropolis-defined accents (safe: these beamercolors exist in metropolis)
\setbeamercolor{progress bar}{fg=deepred}
\setbeamercolor{title separator}{fg=deepred}
% frame number as current/total, in deep red
\metroset{numbering=fraction}
\setbeamercolor{footline}{fg=deepred}
% backup appendix: keep the main frame total at the running-order count
\usepackage{appendixnumberbeamer}
% VISIBLE link colours for rehearsal (internal hyperlinks in deep red)
\hypersetup{colorlinks=true,linkcolor=deepred,citecolor=deepred,urlcolor=deepred}

% ---- [THEORY-TALK] encoding, language, maths, tables, tikz ----------
\usepackage[utf8]{inputenc}
\usepackage[english]{babel}
\usepackage{amsmath}
\usepackage{amssymb}
\usepackage{amsfonts}
\usepackage{booktabs}
\usepackage[official]{eurosym}   % \euro --- this deck quotes euro magnitudes
\usepackage{tikz}
\usetikzlibrary{decorations.pathreplacing,arrows.meta,positioning,fit,backgrounds}

% ---- [THEORY-TALK] the two-agent sketch palette --------------------
% red = agent i / household A ; blue = agent h / household B
\definecolor{agenti}{RGB}{200,30,40}
\definecolor{agenth}{RGB}{20,80,200}

% ---- [THEORY-TALK] notes on a second screen (rehearsal view) -------
% The rehearsal build turns this on; the projection build leaves it off.
\usepackage{pgfpages}
% Long notes overflow the note column by default; render them smaller,
% ragged-right, with tight paragraph spacing so the full text stays
% legible and inside the frame on the second screen.
%
% This deck's notes are longer than the theory talk's -- several run past
% one note page at \scriptsize and produced overfull \vbox warnings.  The
% template below keeps the theory talk's look and adds two things it
% needs: \tiny with a tighter baseline, and an \allowbreak-friendly
% vertical box so a long note breaks across note pages instead of
% overflowing.  Hyphenation is relaxed so no line runs into the margin.
\setbeamertemplate{note page}{%
  \vspace{0.3em}\par
  {\tiny\setlength{\parskip}{0.45em}\setlength{\parindent}{0pt}%
   \setlength{\emergencystretch}{3em}%
   \linespread{1.05}\selectfont
   \raggedright\insertnote\par}%
}
% Notes are prose, not code: let TeX hyphenate freely rather than push a
% long word past the measure.
\hyphenpenalty=50
\tolerance=2000

\DeclareMathOperator*{\argmax}{arg\,max}

% ====================================================================
% ADDITIONS for this deck (not in the theory talk)
% ====================================================================

% ---- channel colours ------------------------------------------------
% One colour per decomposition channel, used consistently on every slide
% that names a channel.  Preferences take the theory talk's agent-red so
% the two decks agree on which side of the argument is which.
\definecolor{chpref}{RGB}{200,30,40}     % preferences   (= agenti)
\definecolor{chacc}{RGB}{20,80,200}      % job access    (= agenth)
\definecolor{chearn}{RGB}{200,120,20}    % earning opportunities
\definecolor{chneeds}{RGB}{40,120,70}    % endowments and needs
\definecolor{chenv}{RGB}{60,60,70}       % the whole environment

% ---- the paper's notation ------------------------------------------
% Audience-facing channel names are WORDS on the slides; these commands
% exist so the words are spelled the same way every time they appear.
\newcommand{\pref}{\textcolor{chpref}{preferences}}
\newcommand{\Pref}{\textcolor{chpref}{Preferences}}
\newcommand{\acc}{\textcolor{chacc}{job access}}
\newcommand{\Acc}{\textcolor{chacc}{Job access}}
\newcommand{\earn}{\textcolor{chearn}{earning opportunities}}
\newcommand{\Earn}{\textcolor{chearn}{Earning opportunities}}
\newcommand{\needs}{\textcolor{chneeds}{household endowments and needs}}
\newcommand{\Needs}{\textcolor{chneeds}{Household endowments and needs}}
\newcommand{\env}{\textcolor{chenv}{the environment}}
\newcommand{\Env}{\textcolor{chenv}{The environment}}

% Symbols.  The theory-paper symbols W, z, R, A, y appear ONLY where the
% companion paper is cited; the estimation notation is separate.
\newcommand{\Wmeasure}{W}                      % theory-paper measure
\newcommand{\Jset}{\mathcal{J}}                % theory-paper job space
\newcommand{\Wone}{W^{1}}                      % theory-paper W^1
\newcommand{\yref}{\mathbf{y}^{w}}             % theory-paper reference pay
\newcommand{\gopp}{g}                          % estimated offer density
\newcommand{\uutil}{u}                         % preference index
\newcommand{\Wmm}{\mathrm{W}}                  % money-metric welfare level

% ---- headline / caveat call-outs -----------------------------------
\newcommand{\headline}[1]{{\large\textcolor{deepred}{#1}}}
\newcommand{\caveat}[1]{{\footnotesize\textcolor{deepred}{#1}}}
% a number with its numerical-integration band
\newcommand{\wband}[2]{#1\,\ensuremath{\pm}\,#2}

% ---- the 25-minute variant ------------------------------------------
% Every frame carries \shortdeck (in the 25-minute running order) or
% \longdeck (45-minute only).  \DeckShort, set by the driver file, drops
% the \longdeck frames; the default keeps everything.
%
% The frozen plan's 25-minute running order is 16 slides, reached from the
% 23 by FIVE cuts and TWO merges:
%   CUT   -> backup : 9, 12, 17, 18, 22          (\longdeck)
%   MERGE           : 6+7 -> one slide, 10+11 -> one slide
% So three tags, not two.  A merged pair is written as two full frames for
% the 45-minute deck; in the short deck the SECOND of the pair drops and
% the first shows its merged content.  23 - 5 - 2 = 16.
%
% Usage:  \shortdeck{\begin{frame}...\end{frame}}     kept in both
%         \longdeck{\begin{frame}...\end{frame}}      45-minute only (cut)
%         \mergedaway{\begin{frame}...\end{frame}}    45-minute only (merged)
%         \ifDeckShort ... \else ... \fi              inline merged content
\newif\ifDeckShort\DeckShortfalse
\newcommand{\shortdeck}[1]{#1}                       % always kept
\newcommand{\longdeck}[1]{\ifDeckShort\else#1\fi}    % cut to backup at 25 min
\newcommand{\mergedaway}[1]{\ifDeckShort\else#1\fi}  % folded into the previous
% content shown only in the merged (25-minute) form of a slide
\newcommand{\mergedonly}[1]{\ifDeckShort#1\fi}

% ---- figure helper ---------------------------------------------------
% Every figure is a PDF from the frozen sprint set, included by name from
% beamer/figures/.  \deckfig{name}{width fraction} keeps the call sites
% uniform and makes a missing file fail loudly at compile time.
\graphicspath{{figures/}}
% Heights leave room for the frametitle, the one caption line every figure
% slide carries, and the metropolis footline -- a taller box overflows the
% frame and shows up as an overfull \vbox.
\newcommand{\deckfig}[2]{%
  \begin{center}\includegraphics[width=#2\textwidth,height=0.70\textheight,%
    keepaspectratio]{#1}\end{center}}
\newcommand{\deckfigtwo}[4]{%
  \begin{columns}[c,onlytextwidth]
  \column{0.49\textwidth}\includegraphics[width=\textwidth,height=0.60\textheight,%
    keepaspectratio]{#1}\\{\centering\scriptsize #2\par}
  \column{0.49\textwidth}\includegraphics[width=\textwidth,height=0.60\textheight,%
    keepaspectratio]{#3}\\{\centering\scriptsize #4\par}
  \end{columns}}

% ---- tables ---------------------------------------------------------
% Several slides carry a comparison table that is naturally wider than the
% text block.  \decktable shrinks one to the measure rather than letting it
% run into the margin (an overfull \hbox), and never enlarges a narrow one.
\newcommand{\decktable}[1]{%
  \begin{center}\resizebox{\textwidth}{!}{#1}\end{center}}

\endinput

% deck_numbers_r6.tex -- GENERATED by make_deck_numbers_r6.py.
% Do not edit by hand.  Every numeral on an R6 slide comes from here.
% Sources: S11 parameter tables (model of record);
%          positive_fit_diagnostics_v2 (G2-ADEQUATE statistics only);
%          welfare/baseline_f1_v1 (BASELINE-F-1, commit 6048c9f, verified b5550af).

% --- BASELINE-F-1 (E1): household EUR/month, unequivalised, dwt-weighted ---
\newcommand{\WFSingN}{1\,540}
\newcommand{\WFSingMean}{1\,435}
\newcommand{\WFSingMedian}{1\,320}
\newcommand{\WFSingGini}{0.238}
\newcommand{\WFSingWorkers}{1\,336}
\newcommand{\WFSingNonworkers}{204}
\newcommand{\WFSingCone}{0}
\newcommand{\WFSingCtwo}{0}
\newcommand{\WFSingCfive}{1.8e-15}
\newcommand{\WFCoupN}{2\,223}
\newcommand{\WFCoupMean}{2\,496}
\newcommand{\WFCoupMedian}{2\,356}
\newcommand{\WFCoupGini}{0.210}
\newcommand{\WFCoupWorkers}{2\,173}
\newcommand{\WFCoupNonworkers}{50}
\newcommand{\WFCoupCone}{0}
\newcommand{\WFCoupCtwo}{0}
\newcommand{\WFCoupCfive}{1.8e-15}

% --- positive fit: ONLY G2-ADEQUATE statistics (weighted, the binding rule) ---
\newcommand{\FitExtCF}{89.7}
\newcommand{\FitExtRatioCF}{0.17}
\newcommand{\FitExtSF}{86.8}
\newcommand{\FitExtRatioSF}{0.24}
\newcommand{\FitAdequateCount}{2}
\newcommand{\FitLimitedCount}{2}

% --- S11 objectives reproduced by the diagnostics evaluator ---
\newcommand{\ObjSingles}{6253.463}
\newcommand{\ObjCouples}{10283.034}
\newcommand{\FitNodes}{100}
\newcommand{\FitSims}{500}
\newcommand{\FitBootReps}{200}
\newcommand{\EvaluatorCommit}{55bb0d0}

% --- S11 model of record: estimated opportunity kernel ---
\newcommand{\SSingParams}{52}
\newcommand{\SSingActive}{41}
\newcommand{\SElevenAccessConst}{-3.174}
\newcommand{\SElevenAccessConstZ}{-7.8}
\newcommand{\SElevenAccessSur}{-1.442}
\newcommand{\SElevenAccessSurZ}{-6.1}
\newcommand{\SElevenHoursFThirtyFive}{2.066}
\newcommand{\SElevenHoursFThirtyFiveZ}{23.1}
\newcommand{\SElevenBetaC}{2.039}
\newcommand{\SElevenBetaCZ}{7.0}
\newcommand{\SElevenWageSigma}{0.382}
\newcommand{\SElevenWageSigmaZ}{28.7}
\newcommand{\SElevenWageEducH}{0.149}
\newcommand{\SElevenWageEducHZ}{4.8}
\newcommand{\SElevenWagePexp}{0.241}
\newcommand{\SElevenWagePexpZ}{2.8}
\newcommand{\SCoupParams}{47}
\newcommand{\SCoupActive}{47}

% --- estimation frames (S11 corrected criterion-A) ---
\newcommand{\FrameSingles}{1\,540}
\newcommand{\FrameCouples}{2\,223}


\graphicspath{{figures/}{figures/slides/}{figures/r6/}}
\ifdefined\RehearsalDeck
  \setbeameroption{show notes on second screen=right}
\fi
\newcommand{\presenternotesize}{\small}
\setbeamertemplate{note page}{%
  \vspace{0.3em}\par
  {\presenternotesize\setlength{\parskip}{0.45em}\setlength{\parindent}{0pt}%
   \setlength{\emergencystretch}{3em}\linespread{1.05}\selectfont
   \raggedright\insertnote\par}}
\setbeamerfont{frametitle}{size=\normalsize,series=\bfseries}
\setbeamertemplate{frametitle}{%
  \nointerlineskip
  \begin{beamercolorbox}[wd=\paperwidth,sep=0pt]{frametitle}%
    \vskip0.65em
    \hspace*{1.2em}\begin{minipage}{0.91\paperwidth}%
      \raggedright\hyphenpenalty=10000\exhyphenpenalty=10000
      \insertframetitle\par
    \end{minipage}\par
    \vskip0.65em
  \end{beamercolorbox}}
\renewcommand{\caveat}[1]{{\fontsize{8.2}{10}\selectfont\raggedright\color{deepred}#1\par}}
\newcommand{\wfunits}{\caveat{Units: household EUR/month, unequivalised;
  survey weight \texttt{dwt}. Descriptive baseline; no counterfactual,
  no decomposition, no causal reading.}}
\newcommand{\wfunitseq}{\caveat{Units: EUR/month, household-equivalised
  (modified-OECD scale); survey weight \texttt{dwt}. A descriptive dispersion
  comparison, not a welfare-loss index. Equivalence scale: the ratified
  modified-OECD scale (Deputy ruling ``SCALE CLOSED; CHILD-SHIFTER FRAMING'',
  s1; \texttt{JMP\_SCALE\_REVIEW\_1\_equivalence\_scale\_economics\_v1.md}).
  Singles and couples equivalised levels are never compared.}}

\title{Unequal Job Opportunities and the Measurement of Welfare Inequality}
\author{Hisham Haydar}
\date{}

\begin{document}

% ---------------------------------------------------------------- title
\begin{frame}[plain]
\vfill
{\LARGE\bfseries\inserttitle\par}
\vspace{1.3em}
{\color{deepred}\rule{0.80\textwidth}{1pt}}\par
\vspace{1.3em}
{\large Hisham Haydar\par}
LISER and University of Luxembourg\par
\vspace{1.3em}
{\small France, EU-SILC priced through EUROMOD.\par}
\vspace{0.6em}
{\footnotesize Seminar of 17 September. Welfare content is presented under
ruling R6; DECOMP-2 is preliminary model-based accounting.\par}
\vfill
\note{Thank you. The question is how much of the inequality in money-metric
well-being is attributable to unequal job opportunities rather than to
differences in taste, and what happens to that attribution when a model
leaves opportunities out. I will show the model, the estimates, the fit, the
welfare definition, a descriptive welfare baseline and the bounded DECOMP-2
accounting exercise. The final decomposition of total well-being inequality
is not ready, and I will say exactly where that boundary lies.}
\end{frame}

% ================================================================ 1 QUESTION
\section{Question and contribution}

\begin{frame}
\headlineframe{Two people with the same tastes and the same wage need not
face the same set of jobs.}
\vfill
\begin{itemize}\itemsep0.7em
  \item How much of measured welfare inequality is attributable to unequal
        \acc{} and \earn{}, rather than to \pref{}?
  \item What does a model that assumes a common choice set do to that
        attribution?
\end{itemize}
\vfill
\caveat{Both questions are about a \emph{measure}, not about a policy
counterfactual.}
\note{The two questions are the same two questions as in every version of
this paper; nothing in the recent rulings touched them. The first is an
attribution question about a money-metric inequality index. The second is a
methodological question: a discrete-choice labour supply model that gives
every household the same alternatives has to explain non-employment with a
taste parameter, and I want to know whether that relabelling survives into
the welfare accounting.}
\end{frame}

\begin{frame}
\headlineframe{The contribution is a structural opportunity set plus a
welfare measure defined on it.}
\vfill
\begin{itemize}\itemsep0.6em
  \item A random-utility, random-opportunity (RURO) model in which
        \emph{which jobs a household can reach} is estimated, not assumed.
  \item The Measure-1 money metric; in this empirical specification its direct
        reference collapses to universally available non-employment.
  \item A bounded three-factor exercise over systematic utility heterogeneity,
        coarse geographic/temporal access and earning opportunities.
\end{itemize}
\vfill
\caveat{The three-factor exercise is preliminary model-based accounting, not
the final decomposition of total well-being inequality.}
\note{The contribution has three parts. Preferences and opportunity
components are jointly estimated, with their separation relying on maintained
functional-form and exclusion restrictions. The welfare measure is taken from
the companion theory paper with Francois Maniquet. Within the bounded DECOMP-2
game, earning-opportunity heterogeneity has a larger Shapley contribution than
the model's coarse geographic/temporal access channel. The final
counterfactual-attainment estimand remains under design.}
\end{frame}

% ================================================================ 2 DATA/MODEL
\section{Data and the structural RURO model}

\begin{frame}
\headlineframe{France, EU-SILC priced through EUROMOD: \FrameSingles{}
single-adult and \FrameCouples{} couple households.}
\vfill
\begin{itemize}\itemsep0.6em
  \item Estimation frames: corrected criterion-A panels, S11 model of record.
  \item Estimation frame: each household's observed job plus \FrameDraws{}
        drawn alternatives, every one priced through the tax--benefit system.
  \item Fit diagnostics use a separate, larger panel: \FitNodes{} common
        quadrature nodes per household.
  \item Weights are the survey weights \texttt{dwt} throughout.
\end{itemize}
\vfill
\caveat{S11 is the single model of record on these slides. Earlier S8 and
R240 frames and their estimates are not shown and are not comparable.}
\note{Two estimation samples, singles and couples, on the corrected
criterion-A frames. Everything on the slides after this point is S11. I am
deliberately not showing the older S8 or R240 numbers alongside, because the
frames differ and a side-by-side would invite a comparison the samples do not
support. Each household contributes its observed job and a hundred drawn
alternatives, each of which is run through EUROMOD so the consumption
argument is a real disposable income, not an approximation. The second
number on the slide is a different object: the positive-fit diagnostics
(POSFIT v2b) integrate over a common panel of two thousand and forty-eight
priced quadrature nodes per household, not over the estimation draws.}
\end{frame}

\begin{frame}
\headlineframe{A job is a package; a household draws from its own
distribution over packages.}
\vfill
\begin{center}
\resizebox{\textwidth}{!}{%
\begin{tikzpicture}[>=Stealth,node distance=9mm,
  box/.style={draw,rounded corners=2pt,align=center,inner sep=5pt,
              font=\footnotesize,minimum height=8mm}]
  \node[box,fill=chacc!10,draw=chacc] (acc) {\textbf{access}\\ employment mass};
  \node[box,fill=chacc!10,draw=chacc,right=of acc] (hrs) {\textbf{hours}\\ piecewise, 35h peak};
  \node[box,fill=chearn!12,draw=chearn,right=of hrs] (occ) {\textbf{occupation}\\ sex-specific};
  \node[box,fill=chearn!12,draw=chearn,right=of occ] (wag) {\textbf{wage offer}\\ log-normal};
  \node[box,fill=black!4,below=14mm of hrs,xshift=10mm,minimum width=52mm] (g)
       {opportunity density $\gopp_i(\cdot)$ over reachable jobs};
  \node[box,fill=chpref!10,draw=chpref,right=32mm of g,minimum width=38mm] (u)
       {preferences $\uutil_i(c,\ell)$};
  \node[box,fill=black!4,below=12mm of g,xshift=26mm,minimum width=68mm] (ch)
       {observed job $=\argmax$ over the drawn set};
  \draw[->] (acc) -- (g); \draw[->] (hrs) -- (g);
  \draw[->] (occ) -- (g); \draw[->] (wag) -- (g);
  \draw[->] (g) -- (ch); \draw[->] (u) -- (ch);
\end{tikzpicture}}
\end{center}
\vfill
\caveat{Access and capability are not separately identified; the access
factor is their product.}
\note{This is the architecture. A job is not an hours cell: it is hours,
occupation and a wage, and whether the household can reach it at all. The
opportunity density factorises into an access margin, a piecewise-constant
hours factor with a peak at the statutory thirty-five hour week, sex-specific
occupation availability, and an occupation-conditioned log-normal wage offer.
Preferences enter separately. The observed job is the argmax over what was
drawn. The identification caveat is on the slide and I keep it there all
talk: the access factor mixes personal capability and market availability and
the cross-section does not separate them.}
\end{frame}

\begin{frame}
\headlineframe{Preferences and the opportunity density are estimated in one
likelihood.}
\slideeq{$\displaystyle
  \Pr(j_i \mid x_i) \;\propto\;
  \exp\!\big\{\, \uutil_i(c_{ij},\ell_{ij}) \;+\; \log \gopp_i(j)
                 \;-\; \log q_i(j) \,\big\}$}
\glossrow{$\uutil$ preferences \quad $\gopp$ estimated opportunity density
\quad $q$ numerical proposal density, divided out}
\vfill
\caveat{The proposal correction is what makes the drawn alternatives a
simulator for the estimated set rather than the choice set itself.}
\note{One likelihood, two structural objects. The proposal density is a
numerical device: the alternatives are drawn from it and it is divided back
out, so it is not part of the economics. That separation matters later,
because the welfare measure I am going to define must not depend on the
proposal density, and I will show that it does not. Preferences and
opportunity components are jointly estimated, with their separation relying
on maintained functional-form and exclusion restrictions. No causal
interpretation follows from that separation. The baseline deliberately keeps
systematic preference heterogeneity parsimonious: age profiles for all four
adult groups and a child-related shifter for women. The child term is
interpreted as a reduced-form behavioural/time-constraint shifter, not as pure
taste. More explicitly, systematic leisure heterogeneity contains an
intercept, age, age squared and group-specific leisure curvature for all four
groups, plus a child shifter for women; sex and household type are carried by
separate parameter blocks.}
\end{frame}

% ================================================================ 3 FIT
\section{Positive estimates and fit}

\begin{frame}
\headlineframe{S11 model of record: \SSingActive{} free parameters for
singles, \SCoupActive{} for couples.}
\vfill
\decktable{\begin{tabular}{lrr}
\toprule
 & singles & couples \\
\midrule
free parameters        & \SSingActive{}  & \SCoupActive{} \\
estimation households  & \FrameSingles{} & \FrameCouples{} \\
objective at $\hat\theta$ & \ObjSingles{} & \ObjCouples{} \\
\bottomrule
\end{tabular}}
\vfill
\caveat{Objectives reproduce exactly under the independent diagnostics
evaluator, package commit \texttt{\EvaluatorCommit{}}.}
\note{The headline fact about the estimates is that they are reproducible:
an evaluator built separately for the diagnostics work recovers both stored
objective values to the last digit shown. That is the precondition for
everything on the next two slides. The parameter counts differ because the
couples specification pins nothing at the bound while the singles
specification pins eleven coefficients by construction.}
\end{frame}

\begin{frame}
\headlineframe{Extensive-margin accuracy, shown only for the groups where it
passes the quadrature-adequacy gate.}
\slidefig{fitext_r6}
\caveat{Values are shown only where the G2 weighted quadrature gate clears
for this statistic; single men are quadrature-limited and withheld. This is
not a group-count result.}
\note{The pre-registered adequacy gate asks whether numerical integration
error is small relative to sampling variation. I show values only where that
specific extensive-accuracy statistic clears the gate and withhold the
single-men value because it is quadrature-limited. This is not a group-count
result. Aggregate participation and occupation margins are reproduced
reasonably closely, while hours fit is uneven. Individual-level predictive
diagnostics raise possible excess stochastic dispersion in some groups; the
group-level classification is not used as a headline result. If asked:
coupled men pass the weighted gate narrowly (ratio
\FitExtRatioCMWeightedVThree{} against 0.25) and miss the unweighted one
(\FitExtRatioCMUnweightedVThree{}); the deck applies the weighted rule
throughout. Source: POSFIT v3 (MNL\_posfit, branch diagnostics/posfit-v3,
commit \texttt{\PosfitVThreeCommit{}}); the extensive-margin numbers are
unchanged from v2b (commit \texttt{\PosfitCommit{}}).}
\end{frame}

\begin{frame}
\headlineframe{Calibration conditioned on prediction: observed hours against
predicted hours, by predicted decile.}
\slidefig{calibration_r6}
\caveat{support coverage audited (POSFIT v2b); calibration statistics partly quadrature-limited; fit verdicts open pending Deputy review}
\note{Households are sorted into deciles of their own predicted expected
hours conditional on working, and each decile's mean observed hours is
plotted against its mean prediction, on the v2b common-node support. The
dashed line is the forty-five degree line. Conditioning on the prediction
rather than on the outcome is the direction that does not mechanically
manufacture regression to the mean. The caption is the honest status. The
support audit has now been done: v2b re-evaluates the fit on a common
quadrature panel far larger than the estimation draws, no household has an
effective sample size below thirty, singles have no unsupported observed
bin, and only a handful of couples fall in a region that finite panel does
not represent. The calibration statistics behind this picture remain partly
quadrature-limited, and the group-level classification is not used as a
headline result. Aggregate participation and occupation margins are reproduced
reasonably closely, while hours fit is uneven. So read the shape, not the gaps.}
\end{frame}

% ================================================================ 4 KERNEL
\section{Evidence for heterogeneous opportunities}

\begin{frame}
\headlineframe{The access kernel is estimated, and it varies across
households.}
\slidefig{kernelacc_r6}
\caveat{S11 singles; horizontal bars are 1.96 $\times$ CR1 cluster-robust
standard errors. Structural and model-conditional; not causal.}
\note{This is the access block of the estimated opportunity density. The
point is not any individual coefficient, it is that the block is estimated
rather than imposed and that it moves: households in different regional
access environments, and in different survey strata, face materially
different employment masses at the same preferences and the same wage. This
is the evidence that the opportunity side of the model is doing work. It is
also the evidence the welfare measure later rests on, which is why I show it
before the measure and not after.}
\end{frame}

\begin{frame}
\headlineframe{So is the earning-opportunity kernel: an offer distribution,
not a single wage.}
\slidefig{kernelwage_r6}
\caveat{S11 singles; occupation terms shift the offer location.
Dispersion $\hat\sigma=\SElevenWageSigma{}$ ($z=\SElevenWageSigmaZ{}$).}
\note{The wage side is an offer density with an estimated dispersion, not a
point wage attached to each person. Education and experience shift the
location, occupation shifts it again, and the residual dispersion is
precisely estimated. Two households with identical observed wages can have
different earning opportunities in this model, because what they face is a
distribution. That is the second half of what I mean by unequal job
opportunities, and together with the access block it is what the welfare
measure has to be defined on.}
\end{frame}

% ================================================================ 5 W1
\section{The welfare measure}

\begin{frame}
\headlineframe{Haydar--Maniquet $\Wone$: the consumption at home that would
leave the household exactly as well off.}
\slideeq{$\Wone_i \;=\; m_i(o;z_i)
  \quad\text{where}\quad
  \uutil_i\big(m_i(o),\,o\big) \;=\; \uutil_i(z_i)$}
\glossrow{$z_i$ attained bundle \quad $o$ the non-employment state \quad
$m_i$ the money metric}
\vfill
\caveat{Mapping F is the primary empirical implementation (ruling R1).
Mapping M is sensitivity only and requires an unidentified intensity
parameter; it is not in the baseline.}
\note{Q\&A note --- ruling R1, JMP\_W1\_fork\_ruling\_v1.md Appendix A.
If asked why this reference and not the market-job reference: the fork
analysis compared Mapping F, the full feasible set including
non-employment, against Mapping M, the market-job universe. R1 rules F
primary for the current paper. M at its dense-law endpoint stops being a
household-specific object and becomes a whole-market reference, and at
finite intensity it needs a set-size parameter the data do not identify.
That is recorded as a limitation and as future work, not as a competing
headline. I am not claiming F is the right reference in the theory; I am
claiming it is the one this empirical specification identifies.}
\end{frame}

\begin{frame}
\headlineframe{Under the current specification $\Wone$ coincides with the
staying-home equivalent --- as a fact about this domain.}
\slideeq{$\Wone_{i,F}
  \;=\; C_i^{\mathrm{obs}}\,
        \exp\!\Big[\tfrac{L_i(j_i^{\mathrm{obs}})-L_i(o)}{\beta_c}\Big]$}
\vfill
\begin{itemize}\itemsep0.4em
  \item Inputs: $C_i^{\mathrm{obs}}$, $L_i(j^{\mathrm{obs}})$, $L_i(o)$,
        $\beta_c$ --- nothing else.
  \item No opportunity density $\gopp$, no proposal $q$, no shock
        $\varepsilon$, no intensity $\kappa$.
  \item Current nonworkers: $W=C$; current workers: $W<C$ under the
        maintained empirical domain.
\end{itemize}
\vfill
\caveat{This is an \textbf{empirical-domain coincidence} under log
consumption and universally available non-employment. It is \textbf{not} the
theorem $\Wone=W^4$ in the Haydar--Maniquet theory.}
\note{Q\&A note --- ruling R1 and BASELINE-F-1 ruling section 2.
The most likely question here is whether I have just proved that measure one
and measure four are the same thing. I have not, and the ruling is explicit
about it. The coincidence has two ingredients, both specific to what is
estimated: consumption enters utility in logs, and non-employment is
behaviourally available to every household. Change either and the two
measures separate. If asked about the consumption argument: ruling E2
confirms that $C^{\mathrm{obs}}$ must be the exact object entering the
estimated utility at the observed state, not a re-chosen normative income
concept, and the measure-correspondence audit closed that provenance before
the baseline was allowed to run. The money metric is derived from the
Measure-1 reference-set principle. Under the current empirical specification
its direct reference collapses to the universally available non-employment
state; opportunity heterogeneity therefore affects the current welfare
measure through attained bundles. For a one-nat shortfall,
$L_i(j_i^{\mathrm{obs}})-L_i(o)=-1$, the illustration is
$W=C_i^{\mathrm{obs}}\exp(-1/\beta_c)<C_i^{\mathrm{obs}}$.}
\end{frame}

% ================================================================ 6 BASELINE
\section{Equivalised $\Wone$-F distribution (primary)}

\begin{frame}
\headlineframe{Baseline $\Wone$-F, equivalised, single adults.}
\vfill
\decktable{\begin{tabular}{lrr}
\toprule
 & $C^{\mathrm{eq}}$ & $\Wone_{F}{}^{\mathrm{eq}}$ \\
\midrule
households        & \WFSingN{}         & \WFSingN{} \\
weighted mean      & \WFSingCEqMean{}   & \WFSingWEqMean{} \\
weighted median    & \WFSingCEqMedian{} & \WFSingWEqMedian{} \\
weighted Gini      & \WFSingCEqGini{}   & \WFSingWEqGini{} \\
\bottomrule
\end{tabular}}
\vfill
\wfunitseq
\caveat{Source: \texttt{docs/results/}
\texttt{JMP\_BASELINE\_F1\_equivalised\_reporting\_v1.md}, commit
\texttt{4c4e07e}, over the verified construction (MNL commit
\texttt{6048c9f}, independently verified \texttt{b5550af}). This is a
descriptive dispersion comparison ($C^{\mathrm{eq}}$ against
$\Wone_F{}^{\mathrm{eq}}$ within the singles sample), not a welfare-loss
index.}
\note{Q\&A note --- BASELINE-F-1 ruling, E3-EQ, and ruling R6.
Equivalised results are primary on this slide: the household-equivalised
Gini of $\Wone_F$ is below the household-equivalised Gini of observed
consumption, for singles. I show that as a descriptive dispersion
comparison only, not as a welfare-loss statement --- the ratio
$\Wone_F/C^{\mathrm{obs}}$ reflects the systematic leisure term at the
observed job relative to home leisure, not a monetary loss, and I am not
constructing a loss index on this slide. The equivalence scale is the
modified-OECD scale, ratified as the primary scale for current JMP
distributional reporting by the Deputy ruling ``SCALE CLOSED; CHILD-SHIFTER
FRAMING'', section 1, on the economics review in
JMP\_SCALE\_REVIEW\_1\_equivalence\_scale\_economics\_v1.md; the scale
question is closed. The unequivalised construction this equivalises is
unchanged and is in the backup appendix.}
\end{frame}

\begin{frame}
\headlineframe{Baseline $\Wone$-F, equivalised, couples.}
\vfill
\decktable{\begin{tabular}{lrr}
\toprule
 & $C^{\mathrm{eq}}$ & $\Wone_{F}{}^{\mathrm{eq}}$ \\
\midrule
households        & \WFCoupN{}         & \WFCoupN{} \\
weighted mean      & \WFCoupCEqMean{}   & \WFCoupWEqMean{} \\
weighted median    & \WFCoupCEqMedian{} & \WFCoupWEqMedian{} \\
weighted Gini      & \WFCoupCEqGini{}   & \WFCoupWEqGini{} \\
\bottomrule
\end{tabular}}
\vfill
\wfunitseq
\caveat{Reported separately from singles: no pooled figure, and no
cross-sample level comparison is drawn on this slide or elsewhere in this
deck. Source: same memo and commits as the singles equivalised slide.}
\note{Q\&A note --- BASELINE-F-1 ruling, E3-EQ, and ruling R6.
Same descriptive comparison, on the couples sample. I am deliberately not
placing this table next to the singles one: the two are separate estimation
frames, the welfare unit differs, and R6 forbids presenting a singles/couples
equivalised-level comparison as a finding. If asked how the couples numbers
compare to singles, the honest answer is that this deck does not draw that
comparison, and I would want to see it examined outside a seminar slide
before treating it as informative. The scale is the same ratified
modified-OECD scale as on the singles slide (Deputy ruling ``SCALE CLOSED;
CHILD-SHIFTER FRAMING'', section 1); ratifying the scale does not license a
cross-sample level comparison.}
\end{frame}

% ================================================================ 7 ARCHITECTURE
\section{Decomposition architecture}

\begin{frame}
\headlineframe{DECOMP-2 equalises three bounded channels, not whole job access.}
\vfill
\begin{itemize}\itemsep0.6em
  \item[$P$] systematic utility heterogeneity (tastes + reduced-form time
        constraints) $\rightarrow$ common reference
  \item[$A$] local geographic/temporal access shifters
        (region, urban/rural, year) $\rightarrow$ common reference
  \item[$B$] earning opportunities $\rightarrow$ common reference offer
        distribution
\end{itemize}
\vfill
\caveat{Household resources, needs and composition are held fixed. Personal
occupation access, hours-band access and node-level alternative
characteristics are not equalised by $A$.}
\note{Q\&A note --- ruling R6 and final economics review M5.
For DECOMP-2, $A$ is local geographic/temporal access shifters (region,
urban/rural, year). It is not personal occupation access, hours-band access
or full job access. Children enter the current structural specification
through a reduced-form female time-constraint shifter. In a unitary
labour-supply model this term may capture both tastes and
childcare/home-production constraints; we therefore do not interpret it as
pure preference heterogeneity. The baseline deliberately keeps systematic
preference heterogeneity parsimonious: age profiles for all four adult groups
and a child-related shifter for women. The child term is interpreted as a
reduced-form behavioural/time-constraint shifter, not as pure taste.}
\end{frame}

\begin{frame}
\headlineframe{DECOMP-2 is executed; the final attainment estimand remains
under design.}
\vfill
\begin{center}
\begin{tikzpicture}[>=Stealth,node distance=7mm,
  b/.style={draw,rounded corners=2pt,align=center,font=\footnotesize,
            inner sep=5pt,minimum height=8mm}]
  \node[b,fill=chenv!8] (s) {coalition $S$};
  \node[b,fill=chenv!8,right=of s] (e) {counterfactual\\environment};
  \node[b,fill=deepred!10,draw=deepred,right=of e] (z) {carried realised\\draws};
  \node[b,fill=chacc!10,right=of z] (w) {$\Wone_F(z_i^{S})$};
  \node[b,fill=chacc!10,right=of w] (i) {inequality $I_S$};
  \draw[->] (s) -- (e); \draw[->] (e) -- (z);
  \draw[->] (z) -- (w); \draw[->] (w) -- (i);
  \node[below=6mm of z,font=\scriptsize,color=deepred,align=center]
       {bounded DECOMP-2 route};
\end{tikzpicture}
\end{center}
\vfill
\caveat{DECOMP-2 reuses the already-priced estimation panel with no
re-estimation and no new pricing. The final decomposition's
counterfactual-attainment estimand remains under design (R5).}
\note{Q\&A note --- ruling R5 and final economics review M10.
The preliminary exercise carries the household's realised draws through each
counterfactual environment. That is an executed, bounded accounting route,
not the final counterfactual-attainment estimand. The final design still has
to choose between a realised-bundle route with conditioned latent
heterogeneity and an ex-ante $g$-computation route; the inequality of expected
welfare is not the expected inequality of welfare.}
\end{frame}

% ================================================================ 8 WIP
\section{Preliminary decomposition and boundary}

\begin{frame}
\headlineframe{Earning opportunities exceed coarse access within the bounded game.}
\vfill
{\small
In a preliminary three-factor structural exercise that holds household
resources, needs and composition fixed, equalising systematic utility
heterogeneity, coarse geographic/temporal access heterogeneity and earning
opportunities changes money-metric well-being inequality by
\DTwoMinPct--\DTwoMaxPct\% of the baseline Gini, depending on household type
and reporting scale. Within the Shapley allocation of that movable component,
earning-opportunity heterogeneity has a larger contribution than the coarse
geographic/temporal access channel in both samples and both reporting
conventions. The preference contribution is not sign-robust to
equivalisation.\par}
\vfill
\caveat{These are preliminary model-based accounting results, not causal
estimates and not the final decomposition of total well-being inequality.}
\note{Q\&A note --- ruling R6 and final economics review M3--M11.
The decomposition is bounded by design because household resources, needs
and composition are held fixed. Within that bounded game, equalising P/A/B
changes the Gini by \DTwoMinPct--\DTwoMaxPct\% of its baseline level, depending
on sample and reporting scale. The restriction is by construction; the
realised numerical magnitude is a model result under that restriction.
Dispersion in consumption is quantitatively large relative to dispersion in
log well-being: Var(log C) is roughly
\DTwoVarMinPct--\DTwoVarMaxPct\% of Var(log W), with the excess offset by a
large negative covariance between consumption and the leisure valuation.
DECOMP-2 holds household resources, needs and composition fixed, leaving
important sources of dispersion outside the P/A/B allocation. If asked about
wage elasticities: a valid gross-wage perturbation requires new tax-benefit
repricing over the affected job alternatives, and the current priced support
does not contain that counterfactual.}
\end{frame}

\begin{frame}
\headlineframe{What I would like from you today.}
\vfill
\begin{itemize}\itemsep0.7em
  \item Is Mapping F the reference you would want for this question?
  \item Which counterfactual attainment estimand would you defend?
  \item What would convince you the access kernel is identified, given that
        access and capability are not separated?
\end{itemize}
\vfill
\caveat{Thank you.}
\note{Q\&A note --- rulings R1, R5, and the identification caveat carried
from the model slide. Three questions, in the order I most need answers. The
reference domain choice is normative and I would like it challenged. The
attainment estimand is the blocking design decision. And the access
identification caveat is the one every discussant raises; I would rather hear
what evidence would settle it than defend the current position.}
\end{frame}

% ================================================================ BACKUP
\appendix

\begin{frame}
\headlineframe{B1 --- Baseline $\Wone$-F, unequivalised, single adults
(verification backup).}
\vfill
\decktable{\begin{tabular}{lr}
\toprule
households                    & \WFSingN{} \\
weighted mean                 & \WFSingMean{} \\
weighted median               & \WFSingMedian{} \\
weighted Gini                 & \WFSingGini{} \\
\midrule
observed workers              & \WFSingWorkers{} \\
observed non-workers          & \WFSingNonworkers{} \\
\bottomrule
\end{tabular}}
\vfill
\wfunits
\caveat{Secondary/verification slide: the equivalised singles slide (\S6) is
primary. Source: \texttt{MNL/outputs/welfare/baseline\_f1\_v1/}
\texttt{baseline\_f1\_full\_sample\_report\_v1.md}, commit
\texttt{6048c9f}, independently verified \texttt{b5550af}.}
\note{This is the literal observed-bundle Mapping-F distribution before
equivalisation, kept here so any question about the un-scaled level or about
worker/non-worker counts has a one-slide answer. It is not the headline: the
equivalised slide in \S6 is. For non-workers the measure equals observed
consumption exactly by construction; for workers it is strictly below it.}
\end{frame}

\begin{frame}
\headlineframe{B2 --- Baseline $\Wone$-F, unequivalised, couples
(verification backup).}
\vfill
\decktable{\begin{tabular}{lr}
\toprule
households                    & \WFCoupN{} \\
weighted mean                 & \WFCoupMean{} \\
weighted median               & \WFCoupMedian{} \\
weighted Gini                 & \WFCoupGini{} \\
\midrule
observed workers              & \WFCoupWorkers{} \\
observed non-workers          & \WFCoupNonworkers{} \\
\bottomrule
\end{tabular}}
\vfill
\wfunits
\caveat{Secondary/verification slide: the equivalised couples slide (\S6) is
primary. Reported separately from singles. No pooled figure: the two samples
are separate estimation frames and the welfare unit differs.}
\note{The couples backup is on its own slide and there is no pooled row:
the two samples are separate estimation frames with different welfare
units, and even on the ratified modified-OECD scale this deck never compares
singles and couples levels. These are unequivalised household totals:
couples pool two earnings streams, and the non-employment count is small.
The pricing-floor issue that affects the neither-work bundle for a subset of
couples does not enter here, because this baseline is evaluated at the
observed bundle and never prices the home state.}
\end{frame}

\begin{frame}
\headlineframe{B3 --- Baseline $\Wone$-F pre-registered checks}
\vfill
\decktable{\begin{tabular}{lll}
\toprule
check & singles & couples \\
\midrule
C1 non-workers, $\max|\Wone-C^{\mathrm{obs}}|$ & \WFSingCone{} & \WFCoupCone{} \\
C2 workers, ratio outside $(0,1)$              & \WFSingCtwo{} & \WFCoupCtwo{} \\
C3 dependency test ($\gopp,q,\varepsilon,\kappa$) & PASS & PASS \\
C4 MEASURE-MAP-1R worked households            & PASS & PASS \\
C5 indifference identity, $\max|\Delta \uutil|$ & \WFSingCfive{} & \WFCoupCfive{} \\
C6 guards triggered                            & 0 & 0 \\
\bottomrule
\end{tabular}}
\vfill
\caveat{All six checks were pre-registered in the BASELINE-F-1 ruling before
execution. C2 violations were to be explained, never clipped; none occurred.}
\note{The checks are the reason the baseline is presentable. C1 and C2 are
the model-implied domain conditions. C3 is the dependency test: the baseline
API accepts no opportunity density, no proposal, no shock and no intensity,
and does not read them indirectly. C5 inverts the utility numerically and
confirms the indifference identity to machine precision.}
\end{frame}

\begin{frame}
\headlineframe{B4 --- The G2 quadrature-adequacy gate}
\vfill
\decktable{\begin{tabular}{lrl}
\toprule
group & extensive accuracy & G2 label \\
\midrule
coupled women & \FitExtCF{}\% & ADEQUATE (ratio \FitExtRatioCF{}) \\
single women  & \FitExtSF{}\% & ADEQUATE (ratio \FitExtRatioSF{}) \\
coupled men   & \FitExtCM{}\% & ADEQUATE (ratio \FitExtRatioCM{}) \\
single men    & ---           & QUADRATURE-LIMITED, withheld \\
\bottomrule
\end{tabular}}
\vfill
\caveat{Weighted, all households. Ratio is Monte Carlo s.e. over sampling
s.d.; the gate is $\le 0.25$. Node bootstrap: \FitBootReps{} replicates over
\FitNodes{} common quadrature nodes per household (diagnostic panel, not the
estimation draws). Simulated bands: \FitSims{} outcome vectors at fixed
$\hat\theta$. POSFIT v2b, commit \texttt{\PosfitCommit{}}.}
\note{The gate was fixed before the statistics were computed. It asks
whether numerical integration error is small enough relative to sampling
variation for the statistic to be a statement about the model rather than
about the quadrature. Withholding rather than downgrading is deliberate:
a quadrature-limited number carries no information about fit.}
\end{frame}

\begin{frame}
\headlineframe{B5 --- Authority for what is and is not on these slides}
\vfill
\decktable{\begin{tabular}{ll}
\toprule
ruling & what it settles here \\
\midrule
R1 & Mapping F primary; $\Wone=W^4$ is a domain coincidence \\
R2 & W1-EA retired from the baseline; preserved, not repaired \\
R4 & MEASURE-MAP-1R: the measure-correspondence gate \\
R5 & counterfactual attainment: in design, not executed \\
R6 & what this seminar may present \\
E1 & BASELINE-F-1 authorised, executed, independently verified \\
E2 & $C^{\mathrm{obs}}$ is the estimated utility's own argument \\
E3 & $C^{\mathrm{obs}}$ benchmark: not authorised, not produced \\
E3-EQ & equivalised reporting (\S6) primary; modified-OECD ratified (SCALE CLOSED, s1) \\
\bottomrule
\end{tabular}}
\vfill
\caveat{R1--R6: \texttt{docs/normative/JMP\_W1\_fork\_ruling\_v1.md},
Appendix A. E1--E3: \texttt{docs/normative/}
\texttt{JMP\_W1\_BASELINE\_F1\_authorization\_and\_Cobs\_ruling\_v1.md}.
E3-EQ: \texttt{docs/results/}
\texttt{JMP\_BASELINE\_F1\_equivalised\_reporting\_v1.md}, commit
\texttt{4c4e07e}.}
\note{This backup slide exists so that any question of the form ``why is
that not here'' has a one-line answer with a document behind it.}
\end{frame}

\end{document}
